\documentclass{article}
\usepackage{../assignment_preamble}

\title{Weekly Problem 1}
\author{Ravi Kini}
\date{April 6, 2026}

\begin{document}

\maketitle

\problem
\subproblem{(a)}
\begin{equation*}
\begin{nd}
\hypo {1} {p \lor q}
\open
\hypo {2} {\lnot p}
\open
\hypo {3} {p}
\have {4} {\bot} \ne{2,3}
\have {5} {q} \be{4}
\close 
\open
\hypo {6} {q}
\have {7} {q} \r{6}
\close 
\have {8} {q} \oe{1,3-5,6-7}
\close
\have {9} {\lnot p \Rightarrow q} \ii{2-8}
\end{nd}
\end{equation*}
We then see that $p \lor q \vdash \lnot p \Rightarrow q$ is a valid argument in classical logic.

\subproblem{(b)}
Let $p = 1$, $q = 0$, and $r = 1$. We see that $q \lor r = 1$, so $p \Rightarrow (q \lor r) = 1$. However, $\lnot p = 0$. Consequently, the argument $p \Rightarrow (q \lor r), \lnot q \vdash \lnot p$ is an invalid argument in classical logic.

\end{document}